% GENERATED FILE. Edit canonical lessons, then run npm run generate:books.
% canonical-source-hash: fnv1a64:96c272438ea55a2e
% canonical-lessons: RU-C368-blagotvoritelnost, RU-C368-volontyor, RU-C368-posobie, RU-C368-zhertvovat, RU-C368-nuzhdatsya, RU-R368-review-361-368

\chapter{Helping people in need}
\label{ch:ru-helping-people-in-need}
\hlchaptermodality{368}

\begin{chapteropening}
\textbf{By the end of this chapter:} \emph{I can talk about charity and volunteers, a state benefit, giving money and people who are in need.}

Complete the last lesson of chapter 368: I can talk about charity and volunteers, a state benefit, giving money and people who are in need.
\end{chapteropening}

\section[\texorpdfstring{\ru{благотворительность} (blagotvorítelnost)}{благотворительность (blagotvorítelnost)}]{\ru{благотворительность} (blagotvorítelnost) — charity}

\label{lesson:RU-C368-blagotvoritelnost}



\begin{quote}
Before the new one: say the Russian for wealth, then the Russian for unemployment.
\end{quote}

\subsection*{You'll want to know: \ru{благотворительность}}
\textbf{\ru{благотворительность}} (\emph{blagotvorítelnost}) — \textquotedblleft{}charity\textquotedblright{}.

\begin{quote}
\ru{Она} \ru{много} \ru{лет} \ru{занимается} \ru{благотворительностью}.
\end{quote}

\begin{quote}
\emph{Oná mnógo let zanimáyetsya blagotvorítelnostyu.}
\end{quote}

\begin{quote}
She has done charity work for many years.
\end{quote}

Feminine. \emph{\ru{Благотворительный} \ru{концерт}} — a charity concert.

\begin{cousinweb}
\textbf{\ru{благотворительность}} is built from \emph{\ru{благо}} \textquotedblleft{}good\textquotedblright{}, a Church Slavonic word with \emph{-\ru{ла}-}, and \emph{\ru{творить}} \textquotedblleft{}to make, to do\textquotedblright{}: doing good.
\end{cousinweb}

\begin{grammarlens}[title={What you've built}]
One more word for telling, explaining and describing at B1.
\end{grammarlens}

\subsection*{Your turn}
\begin{itemize}
\raggedright
  \item \emph{Say it:} \emph{blagotvorítelnost}
  \item \emph{Say it:} \emph{blagotvorítelnost}, once more
  \item \emph{Say it:} a sentence of your own with \emph{blagotvorítelnost}
  \item \emph{Recall:} say the Russian for wealth, then the Russian for unemployment, then say \emph{blagotvorítelnost} again
\end{itemize}

\begin{tcolorbox}[breakable,colback=teal!4,colframe=teal!35!black,title={Before you move on}]
What does \ru{благотворительность} mean? (\textquotedblleft{}Charity\textquotedblright{}.) Say it once more.
\end{tcolorbox}
\section[\texorpdfstring{\ru{волонтёр} (volontyór)}{волонтёр (volontyór)}]{\ru{волонтёр} (volontyór) — a volunteer}

\label{lesson:RU-C368-volontyor}



\begin{quote}
Before the new one: say the Russian for unemployment, then the Russian for charity.
\end{quote}

\subsection*{You'll want to know: \ru{волонтёр}}
\textbf{\ru{волонтёр}} (\emph{volontyór}) — \textquotedblleft{}a volunteer\textquotedblright{}.

\begin{quote}
\ru{Волонтёры} \ru{помогают} \ru{в} \ru{больнице}.
\end{quote}

\begin{quote}
\emph{Volontyóry pomogáyut v bolnítse.}
\end{quote}

\begin{quote}
Volunteers help at the hospital.
\end{quote}

Masculine. \emph{\ru{Доброволец}} means the same.

\begin{cousinweb}
\textbf{\ru{волонтёр}} was borrowed from French \emph{volontaire} (EtymDB) \textquotedblleft{}voluntary, a volunteer\textquotedblright{}, from Latin \emph{voluntās} \textquotedblleft{}will, choice\textquotedblright{} (Lewis \& Short).
\end{cousinweb}

\begin{grammarlens}[title={What you've built}]
One more word for telling, explaining and describing at B1.
\end{grammarlens}

\subsection*{Your turn}
\begin{itemize}
\raggedright
  \item \emph{Say it:} \emph{volontyór}
  \item \emph{Say it:} \emph{volontyór}, once more
  \item \emph{Say it:} a sentence of your own with \emph{volontyór}
  \item \emph{Recall:} say the Russian for unemployment, then the Russian for charity, then say \emph{volontyór} again
\end{itemize}

\begin{tcolorbox}[breakable,colback=teal!4,colframe=teal!35!black,title={Before you move on}]
What does \ru{волонтёр} mean? (\textquotedblleft{}A volunteer\textquotedblright{}.) Say it once more.
\end{tcolorbox}
\section[\texorpdfstring{\ru{пособие} (posóbiye)}{пособие (posóbiye)}]{\ru{пособие} (posóbiye) — a benefit, an allowance; a study guide}

\label{lesson:RU-C368-posobie}



\begin{quote}
Before the new one: say the Russian for charity, then the Russian for a volunteer.
\end{quote}

\subsection*{You'll want to know: \ru{пособие}}
\textbf{\ru{пособие}} (\emph{posóbiye}) — \textquotedblleft{}a benefit, an allowance; a study guide\textquotedblright{}.

\begin{quote}
\ru{Он} \ru{получает} \ru{пособие} \ru{по} \ru{безработице}.
\end{quote}

\begin{quote}
\emph{On polucháyet posóbiye po bezrabótitse.}
\end{quote}

\begin{quote}
He receives unemployment benefit.
\end{quote}

A neuter noun. \emph{\ru{Пособие} \ru{по} \ru{безработице}} — unemployment benefit. \emph{\ru{Учебное} \ru{пособие}} — a study guide.

\begin{cousinweb}
\textbf{\ru{пособие}} is built on the old verb \emph{\ru{пособить}} \textquotedblleft{}to help\textquotedblright{}. The origin of the root is not settled.
\end{cousinweb}

\begin{grammarlens}[title={What you've built}]
One more word for telling, explaining and describing at B1.
\end{grammarlens}

\subsection*{Your turn}
\begin{itemize}
\raggedright
  \item \emph{Say it:} \emph{posóbiye}
  \item \emph{Say it:} \emph{posóbiye}, once more
  \item \emph{Say it:} a sentence of your own with \emph{posóbiye}
  \item \emph{Recall:} say the Russian for charity, then the Russian for a volunteer, then say \emph{posóbiye} again
\end{itemize}

\begin{tcolorbox}[breakable,colback=teal!4,colframe=teal!35!black,title={Before you move on}]
What does \ru{пособие} mean? (\textquotedblleft{}A benefit, an allowance; a study guide\textquotedblright{}.) Say it once more.
\end{tcolorbox}
\section[\texorpdfstring{\ru{жертвовать} (zhértvovat)}{жертвовать (zhértvovat)}]{\ru{жертвовать} (zhértvovat) — to donate, to give; to sacrifice}

\label{lesson:RU-C368-zhertvovat}



\begin{quote}
Before the new one: say the Russian for a volunteer, then the Russian for a benefit.
\end{quote}

\subsection*{You'll want to know: \ru{жертвовать}}
\textbf{\ru{жертвовать}} (\emph{zhértvovat}) — \textquotedblleft{}to donate, to give; to sacrifice\textquotedblright{}.

\begin{quote}
\ru{Каждый} \ru{год} \ru{мы} \ru{жертвуем} \ru{деньги} \ru{детскому} \ru{дому}.
\end{quote}

\begin{quote}
\emph{Kázhdyy god my zhértvuyem déngi détskomu dómu.}
\end{quote}

\begin{quote}
Every year we donate money to a children’s home.
\end{quote}

Imperfective; the present is \emph{\ru{жертвую}, \ru{жертвуешь}}; the perfective is \emph{\ru{пожертвовать}}. With the instrumental it means to sacrifice: \emph{\ru{жертвовать} \ru{временем}}.

\begin{cousinweb}
\textbf{\ru{жертвовать}} is built on \emph{\ru{жертва}} \textquotedblleft{}a sacrifice, a victim\textquotedblright{}, which is usually traced to an old Slavic verb meaning to praise or to offer to the gods.
\end{cousinweb}

\begin{grammarlens}[title={What you've built}]
One more word for telling, explaining and describing at B1.
\end{grammarlens}

\subsection*{Your turn}
\begin{itemize}
\raggedright
  \item \emph{Say it:} \emph{zhértvovat}
  \item \emph{Say it:} \emph{zhértvovat}, once more
  \item \emph{Say it:} a sentence of your own with \emph{zhértvovat}
  \item \emph{Recall:} say the Russian for a volunteer, then the Russian for a benefit, then say \emph{zhértvovat} again
\end{itemize}

\begin{tcolorbox}[breakable,colback=teal!4,colframe=teal!35!black,title={Before you move on}]
What does \ru{жертвовать} mean? (\textquotedblleft{}To donate, to give; to sacrifice\textquotedblright{}.) Say it once more.
\end{tcolorbox}
\section[\texorpdfstring{\ru{нуждаться} (nuzhdátsya)}{нуждаться (nuzhdátsya)}]{\ru{нуждаться} (nuzhdátsya) — to need, to be in need}

\label{lesson:RU-C368-nuzhdatsya}



\begin{quote}
Before the new one: say the Russian for a benefit, then the Russian for to donate.
\end{quote}

\subsection*{You'll want to know: \ru{нуждаться}}
\textbf{\ru{нуждаться}} (\emph{nuzhdátsya}) — \textquotedblleft{}to need, to be in need\textquotedblright{}.

\begin{quote}
\ru{Эта} \ru{семья} \ru{нуждается} \ru{в} \ru{помощи}.
\end{quote}

\begin{quote}
\emph{Éta semyá nuzhdáyetsya v pómoshchi.}
\end{quote}

\begin{quote}
This family is in need of help.
\end{quote}

Imperfective only. \emph{\ru{Нуждаться} \ru{в}} + prepositional. More formal than \emph{\ru{нужен}}.

\begin{cousinweb}
\textbf{\ru{нуждаться}} is built on \emph{\ru{нужда}} \textquotedblleft{}need, want\textquotedblright{}, the noun behind \emph{\ru{нужный}} \textquotedblleft{}necessary\textquotedblright{}, inherited from Proto-Slavic.
\end{cousinweb}

\begin{grammarlens}[title={What you've built}]
One more word for telling, explaining and describing at B1.
\end{grammarlens}

\subsection*{Your turn}
\begin{itemize}
\raggedright
  \item \emph{Say it:} \emph{nuzhdátsya}
  \item \emph{Say it:} \emph{nuzhdátsya}, once more
  \item \emph{Say it:} a sentence of your own with \emph{nuzhdátsya}
  \item \emph{Recall:} say the Russian for a benefit, then the Russian for to donate, then say \emph{nuzhdátsya} again
\end{itemize}

\begin{tcolorbox}[breakable,colback=teal!4,colframe=teal!35!black,title={Before you move on}]
What does \ru{нуждаться} mean? (\textquotedblleft{}To need, to be in need\textquotedblright{}.) Say it once more.
\end{tcolorbox}
\section[(dialogue)]{Words from chapters 361–368 — a review}

\label{lesson:RU-R368-review-361-368}



\begin{quote}
Say the Russian for an enterprise and profit.
\end{quote}

\subsection*{Your turn}
Say these aloud:

\begin{itemize}
\raggedright
  \item an enterprise, profit, a loss, income, a deal
  \item a promotion, a bonus, a skill, a specialist, to dismiss
  \item negotiations, a contract, a client, a competitor, to cooperate
  \item society, the population, public, social, an inhabitant
  \item the state, the government, power, politics, a politician
  \item an election, to vote, a party, a protest, to protest
  \item freedom, equality, poverty, wealth, unemployment
  \item charity, a volunteer, a benefit, to donate, to need
\end{itemize}

\begin{tcolorbox}[breakable,colback=teal!4,colframe=teal!35!black,title={Before you move on}]
An enterprise? (\textbf{\ru{предприятие}}.) Profit? (\textbf{\ru{прибыль}}.) A loss? (\textbf{\ru{убыток}}.) Income? (\textbf{\ru{доход}}.) A deal? (\textbf{\ru{сделка}}.) A promotion? (\textbf{\ru{повышение}}.) A bonus? (\textbf{\ru{премия}}.) A skill? (\textbf{\ru{навык}}.) A specialist? (\textbf{\ru{специалист}}.) To dismiss? (\textbf{\ru{увольнять}}.) Negotiations? (\textbf{\ru{переговоры}}.) A contract? (\textbf{\ru{контракт}}.) A client? (\textbf{\ru{заказчик}}.) A competitor? (\textbf{\ru{конкурент}}.) To cooperate? (\textbf{\ru{сотрудничать}}.) Society? (\textbf{\ru{общество}}.) The population? (\textbf{\ru{население}}.) Public? (\textbf{\ru{общественный}}.) Social? (\textbf{\ru{социальный}}.) An inhabitant? (\textbf{\ru{житель}}.) The state? (\textbf{\ru{государство}}.) The government? (\textbf{\ru{правительство}}.) Power? (\textbf{\ru{власть}}.) Politics? (\textbf{\ru{политика}}.) A politician? (\textbf{\ru{политик}}.) An election? (\textbf{\ru{выборы}}.) To vote? (\textbf{\ru{голосовать}}.) A party? (\textbf{\ru{партия}}.) A protest? (\textbf{\ru{протест}}.) To protest? (\textbf{\ru{протестовать}}.) Freedom? (\textbf{\ru{свобода}}.) Equality? (\textbf{\ru{равенство}}.) Poverty? (\textbf{\ru{бедность}}.) Wealth? (\textbf{\ru{богатство}}.) Unemployment? (\textbf{\ru{безработица}}.) Charity? (\textbf{\ru{благотворительность}}.) A volunteer? (\textbf{\ru{волонтёр}}.) A benefit? (\textbf{\ru{пособие}}.) To donate? (\textbf{\ru{жертвовать}}.) To need? (\textbf{\ru{нуждаться}}.)
\end{tcolorbox}
